%% sections/02d_order_and_shadow.tex
%% Sezione 2.x — Order Analysis and Shadow Hamiltonian
%% Da inserire dopo 02c_time_reversal_proof.tex

\subsection{Order of the Method and Shadow Hamiltonian}
\label{sec:order_and_shadow}
\label{sec:shadow}

\subsubsection{Local Truncation Error via BCH}
\label{sec:bch_analysis}

The Baker-Campbell-Hausdorff (BCH) formula gives the effective
Hamiltonian $\tilde{H}_h$ such that the one-step map
$\Phi_h = \exp(h\,L_{\tilde{H}_h})$ exactly, where
$L_H$ denotes the Lie derivative associated with $H$.
For the leapfrog (second-order Störmer-Verlet) map:
\begin{equation}
  \tilde{H}^{(2)}_h = H + \frac{h^2}{24}
    \bigl\{V,\{V,T\}\bigr\}
    + \frac{h^2}{12}\bigl\{T,\{T,V\}\bigr\}
    + \mathcal{O}(h^4),
  \label{eq:shadow_leapfrog}
\end{equation}
where $\{A,B\} = \nabla_\vp A \cdot \nabla_\vq B
- \nabla_\vq A \cdot \nabla_\vp B$ is the Poisson bracket
\citep{HairerLubichWanner2006}.

For the Yoshida-4 composition with coefficients $w_0, w_1$
(Equation~\ref{eq:yoshida_coeffs}), the BCH expansion gives:
\begin{equation}
  \tilde{H}^{(4)}_h = H + h^4\,\sigma_4\,E_4 + \mathcal{O}(h^6),
  \label{eq:shadow_yoshida}
\end{equation}
where
\begin{equation}
  \sigma_4 = \frac{2w_1^4 + w_0^4}{24}
  \approx 0.041,
  \label{eq:sigma4}
\end{equation}
and $E_4$ is a fourth-order error term involving fourth-order
Poisson brackets of $T$ and $V$. The specific form relevant for
the gravitational potential $V = -\mu/r$ is:
\begin{equation}
  E_4 = \frac{1}{2}\bigl\{T,\{T,\{T,V\}\}\bigr\}
       + \frac{1}{6}\bigl\{V,\{V,\{V,T\}\}\bigr\}.
  \label{eq:E4_brackets}
\end{equation}

\subsubsection{Explicit Form of the Shadow Hamiltonian}
\label{sec:shadow_explicit}

Evaluating the Poisson brackets in \eqref{eq:E4_brackets} for
$T = |\vp|^2/2$ and $V = -\mu/r$:

\paragraph{First bracket.}
$\{T,V\} = \vp \cdot \nabla_\vq V = \vp \cdot (\mu\vq/r^3)
= \mu(\vp\cdot\hat{\vq})/r^2$.

\paragraph{Second bracket.}
$\{T,\{T,V\}\} = \nabla_\vp[\{T,V\}] \cdot \nabla_\vq T$
requires the full computation; the result involves
$\mu|\vp|^2/r^3 - 3\mu(\vp\cdot\hat{\vq})^2/r^3$.

\paragraph{Leading shadow correction.}
After evaluation, the shadow Hamiltonian takes the form:
\begin{equation}
  \boxed{
  \tilde{H} = H + h^4\,\sigma_4
  \Bigl[
    |\nabla\Phi|^2\,\Phi'' + \vp^T\,\mathbf{H}(\Phi)\,\vp\,\Phi''
  \Bigr]
  + \mathcal{O}(h^6),
  }
  \label{eq:shadow_hamiltonian}
\end{equation}
where $\Phi = -\mu/r$ is the gravitational potential,
$\Phi'' = \mu/r^3$,
$|\nabla\Phi|^2 = (\mu/r^2)^2$, and
$\mathbf{H}(\Phi)$ is the Hessian matrix of $\Phi$
(Equation~\ref{eq:A_matrix}).

The shadow Hamiltonian $\tilde{H}$ is exactly conserved by the
discrete flow of $\Phi_h$ (to the order shown), so
$|\Delta\tilde{H}/\tilde{H}_0|$ measures the residual numerical
error after accounting for the leading shadow correction.
The diagnostic
\begin{equation}
  |\Delta\tilde{H}/\tilde{H}_0| \ll |\Delta H/H_0|
  \label{eq:shadow_diagnostic}
\end{equation}
confirms that the integrator is operating in the shadow
Hamiltonian regime, i.e., the physical energy error is
dominated by the $h^4$ correction rather than by
higher-order numerical noise
\citep{HairerLubichWanner2006,Zeebe2023}.

\subsubsection{Fourth-Order Convergence}
\label{sec:fourth_order}

The global error of the Yoshida-4 method satisfies
\begin{equation}
  \|\vq_N - \vq(t_N)\| + \|\vp_N - \vp(t_N)\|
  = \mathcal{O}(h^4),
  \label{eq:global_error}
\end{equation}
uniformly for $t_N = Nh \leq T$ on any finite time interval.
This is a consequence of the shadow Hamiltonian analysis: since
$\tilde{H}^{(4)} - H = \mathcal{O}(h^4)$, the discrete trajectory
lies exponentially close to the exact trajectory of $\tilde{H}$,
which in turn is $\mathcal{O}(h^4)$ close to the exact trajectory
of $H$ \citep[Theorem~IX.3.5]{HairerLubichWanner2006}.

For the adaptive case with step $h = g(\vq)\Delta s$, the
local error per pseudo-time step is
$\mathcal{O}(\Delta s^5)$, giving a global error
\begin{equation}
  \mathcal{O}\!\left(\frac{g_{\max}^4\,\Delta s^4}{S}\right)
  = \mathcal{O}(\varepsilon^4),
  \label{eq:adaptive_global_error}
\end{equation}
where $S$ is the total pseudo-time and we used
$\Delta s \sim \varepsilon/\sqrt{\rho_{\max}(\mathbf{H})}$.
The precision parameter $\varepsilon$ therefore directly controls
the global error, as confirmed empirically in
Figure~\ref{fig:energy_vs_precision}.
